\section{Extreme amenability}\label{sec:ea}

The simplest possible universal minimal flow is a single point. Groups with
this property are far from locally compact, and they arise from concentration
of measure and from structural Ramsey theory.

\begin{definition}[Extremely amenable group]\label{def:ea}
\uses{def:flow}
A topological group $G$ is \emph{extremely amenable} if every $G$-flow has a
fixed point, that is, a point $x$ with $gx=x$ for all $g\in G$. This is also
called the \emph{fixed point on compacta property}.
\end{definition}

\begin{proposition}\label{prop:ea-trivial-umf}
\uses{def:ea, thm:ellis, prop:minimal-subflow}
A topological group $G$ is extremely amenable if and only if $M(G)$ is a
single point.
\end{proposition}

\begin{proof}
If $G$ is extremely amenable, a fixed point of $M(G)$ is a subflow, so
$M(G)$ is a point by minimality. Conversely, if $M(G)$ is a point, let $X$ be
a $G$-flow. By Proposition~\ref{prop:minimal-subflow} it contains a minimal
subflow $Y$, which is a factor of $M(G)$ and hence a single fixed point.
\end{proof}

\begin{remark}\label{rem:ea-amenable}
\uses{def:ea}
A topological group is \emph{amenable} if every $G$-flow carries a
$G$-invariant regular Borel probability measure. Extremely amenable groups are
amenable: the Dirac measure at a fixed point is invariant.
\end{remark}

\begin{remark}[The first examples]\label{rem:herer-christensen}
\uses{def:ea, rem:ea-amenable}
Historically, the first extremely amenable groups were the groups
$L^0(\varphi,\mathbb{R})$ of Herer and Christensen \cite{HererChristensen1975},
built from a pathological submeasure $\varphi$. They have no nontrivial
continuous unitary representations; being abelian, they are amenable, and
together the two properties make them extremely amenable. Natural examples
came later: see Sections~\ref{sec:gromov-milman}, \ref{sec:pestov}
and~\ref{sec:kpt}, and the monograph \cite{Pestov2006}.
\end{remark}
